\let\negmedspace\undefined
\let\negthickspace\undefined
\documentclass[journal]{IEEEtran}
\usepackage[a5paper, margin=10mm, onecolumn]{geometry}
\usepackage{lmodern} % Ensure lmodern is loaded for pdflatex
\usepackage{tfrupee} % Include tfrupee package

\setlength{\headheight}{1cm} % Set the height of the header box
\setlength{\headsep}{0mm}      % Set the distance between the header box and the top of the text

\usepackage{gvv-book}
\usepackage{gvv}
\usepackage{cite}
\usepackage{amsmath,amssymb,amsfonts,amsthm}
\usepackage{algorithmic}
\usepackage{graphicx}
\graphicspath{{./figs/}}
\usepackage{textcomp}
\usepackage{xcolor}
\usepackage{txfonts}
\usepackage{listings}
\usepackage{enumitem}
\usepackage{mathtools}
\usepackage{gensymb}
\usepackage{comment}
\usepackage[breaklinks=true]{hyperref}
\usepackage{tkz-euclide}
\usepackage{listings}
\usepackage{gvv}

\def\inputGnumericTable{}

\usepackage[latin1]{inputenc}

\usepackage{color}

\usepackage{array}

\usepackage{longtable}

\usepackage{calc}
\usepackage{multirow}

\usepackage{hhline}

\usepackage{ifthen}

\usepackage{lscape}
\usepackage{circuitikz}

\begin{document}
	
	\bibliographystyle{IEEEtran}
	\vspace{3cm}
	
	\title{5.13.67}
	\author{EE25BTECH11042 - Nipun Dasari}
	\maketitle
	
	\renewcommand{\thefigure}{\theenumi}
	\renewcommand{\thetable}{\theenumi}
	\setlength{\intextsep}{10pt} % Space between text and floats
	
	
	\numberwithin{equation}{enumi}
	\numberwithin{figure}{enumi}
	\renewcommand{\thetable}{\theenumi}
	
	\textbf{Question}:\\
	For what value of k do the following system of equations possess a non trivial (i.e.,
	not all zero) solution over the set of rationals Q?
	\begin{align}
		x + ky + 3z = 0\\
	3x+ky-2z=0\\
		2x+3y-4z=0\\
	\end{align}
	\solution
	Given:
	\begin{align}
		\begin{myvec}{1\\ k \\3}\end{myvec}^\top\vec{x} = 0 \label{1}
	\end{align}
	\begin{align}
		\begin{myvec}{3\\ k \\-2}\end{myvec}^\top\vec{x} = 0\label{2}
	\end{align}
	\begin{align}
		\begin{myvec}{2\\3\\-4}\end{myvec}^\top\vec{x} = 0\label{3}
	\end{align}
	Combining \eqref{1},\eqref{2},\eqref{3}:
	\begin{align}
		A\vec{x}=0
	\end{align}
	where A is :
	\begin{align}
		\begin{myvec}{1&k&3\\3&k&-2\\2&3&-4}\end{myvec}
	\end{align}
	For this system to have non-trivial solution these planes must intersect on a line, i.e, any two of the rows must be linearly dependent in order to have infinitely many non-trivial solutions. We simplify the matrix by using gaussian elimination:
	\begin{align}
		\begin{myvec}{1&k&3\\3&k&-2\\2&3&-4}\end{myvec}
		&\xleftrightarrow{,R_2 \gets R_2-3R_1}
		\begin{myvec}{1&k&3\\0&-2k&-11\\2&3&-4}\end{myvec}
		&\xleftrightarrow{,R_3 \gets R_3-2R_1}
		\begin{myvec}{1&k&3\\0&-2k&-11\\0&3-2k&-10}\end{myvec}
	\end{align}
	Now, we continue the row reduction to get the matrix into row echelon form:
	\begin{align}
		\begin{myvec}{1&k&3\\0&-2k&-11\\0&3-2k&-10}\end{myvec}
		&\xleftrightarrow{,R_3 \gets 2kR_3+(3-2k)R_2}
		\begin{myvec}{1&k&3\\0&-2k&-11\\0&0&2k-33}\end{myvec}
	\end{align}
	For a non-trivial solution, the third element of the last row must be zero.
	\begin{align}
		2k - 33 = 0
	\end{align}
	\begin{align}
		\implies k = \frac{33}{2}
	\end{align}
	\begin{figure}[H]
		\centering
		\includegraphics[width = 0.9\columnwidth]{figs/Figure_2.png}
		\caption*{}
		\label{fig2}
	\end{figure}
	\begin{figure}[H]
		\centering
		\includegraphics[width = 0.9\columnwidth]{figs/Figure_1.png}
		\caption*{}
		\label{fig1}
	\end{figure}
\end{document}